%
% Based on https://github.com/BaseMax/developer-cv-template (Max Base)
% Tailored for: iFood, Staff Backend Software Engineer (Ref. 100221)
%
\documentclass[a4paper,10pt]{article}
\usepackage[utf8]{inputenc}
\usepackage[T1]{fontenc}
\usepackage{lmodern}
\input{glyphtounicode}
\pdfgentounicode=1
\usepackage{geometry}
\geometry{top=1cm, bottom=1cm, left=0.8cm, right=0.8cm}
\usepackage{parskip}
\usepackage{hyperref}
\hypersetup{
	colorlinks=true,
	urlcolor=blue,
	linkcolor=black
}
\usepackage{enumitem}
\usepackage{needspace}
\setlist[itemize]{noitemsep, topsep=0pt}

\pagestyle{empty}

\newcommand{\sectiontitle}[1]{\needspace{5\baselineskip}\vspace{6pt}\noindent\textbf{\large #1}\vspace{3pt}\hrule\vspace{6pt}}

\newenvironment{smallprojects}{
	\begingroup
	\small
	\setlength{\parskip}{3pt}
}{
	\endgroup
}

\begin{document}

	%=== HEADER ===
	\begin{center}
		{\LARGE \textbf{Rodrigo Monteiro}}\\
		\vspace{2pt}
		São Mateus, ES, Brazil \quad | \quad \href{mailto:rodrigo.monteiro.tech@gmail.com}{rodrigo.monteiro.tech@gmail.com} \quad | \quad \href{https://givemeurdata.uk}{givemeurdata.uk} \\
		\href{https://github.com/TsukiGva2}{github.com/TsukiGva2} \quad | \quad
		\href{https://www.linkedin.com/in/rodrigo-monteiro-0561992b7/}{linkedin.com/in/rodrigo-monteiro-0561992b7} \\
		\vspace{2pt}
		\textit{Backend engineer: Go, distributed edge systems, PostgreSQL, Linux}
	\end{center}

	%=== PROFESSIONAL SUMMARY ===
	Backend engineer who owned a production distributed system end to end. For two years I was the main engineer behind \textbf{MonoTempo}, the RFID race-timing platform behind live sporting events: RFID readers at the start and finish line, edge units in the field, message-driven services in \textbf{Go}, and the cloud API that publishes the rankings. I designed it offline-first because races happen where connectivity is unreliable, then replaced it with a simpler MQTT and cloud-processing design that ended incidents at events. I also built \textbf{PvAPI}, a Go marketplace backend with transactional checkout, row-level locking and a spec-driven AI-agent workflow. Today I work on data and operational intelligence for a university, writing SQL against a large academic database and using \mbox{AI-assisted} development daily. I am comfortable going from firmware to cloud, owning decisions and explaining their trade-offs.

	%=== CORE SKILLS & EXPERTISE ===
	\sectiontitle{Core Skills \& Expertise}
	\begin{itemize}
		\item \textbf{Go}: concurrent services (goroutines, channels), Gin, pgx, message consumers and producers, CLI tools
		\item \textbf{Distributed \& Edge Systems}: offline-first and store-and-forward design, RabbitMQ, MQTT, idempotent re-sending, crash-and-restart supervision
		\item \textbf{Data}: PostgreSQL, MySQL, SQLite (sharded for write throughput), complex SQL for analytics and data-quality checks
		\item \textbf{Infrastructure}: Linux, Docker, Podman (Quadlets), systemd, nginx, shell automation, Ansible (basic)
		\item \textbf{Other languages}: Python, C, C++ (CPython extensions), Lua (OpenResty/Lapis), PHP (Laravel, CodeIgniter)
		\item \textbf{Embedded}: Arduino/AVR, serial protocols with checksums, RFID readers (Chafon, Impinj), nanoFORTH
		\item \textbf{AI}: designed a spec-driven workflow for coding agents (specs as issues, test-gated guardrails, custom skills); daily \mbox{AI-assisted} development; deployed a local text-to-speech model on edge units
	\end{itemize}

	%=== EXPERIENCE ===
	\sectiontitle{Relevant Experience}

	\textbf{Systems Operator -- Centro Universitário Vale do Cricaré (UNIVC), Hub do Cricaré} \hfill São Mateus, ES (on-site)\\
	\textit{Jun 2026 -- present}
	\begin{itemize}
		\item Part of the team that turns raw data from the university's academic platforms into information for management, faculty and students.
		\item Built the student-data extraction API over the academic ERP (SEI) database: asynchronous CSV exports behind token auth, which replaced a web scraper as the data source for the team's enrollment analytics dashboard.
		\item Write and optimize SQL queries on the academic system (SEI) to extract operational metrics, track enrollment, validate classes and check teaching schedules.
		\item Build analytical reports that monitor use of the SEI and the virtual learning environment (AVA), supporting bug diagnosis, data consistency and content availability.
		\item Investigate level 1 and 2 incidents through system logs and tables, finding the root cause of access and navigation problems with the support team.
		\item Map recurring failures and inconsistent records with SQL, enabling preventive fixes in the database that reduce ticket volume.
	\end{itemize}

	\textbf{Software Engineer, Backend \& IoT -- MyTempo (RFID race timing)} \hfill São Mateus, ES\\
	\textit{2024 -- 2026}
	\begin{itemize}
		\item Main engineer of \textbf{MonoTempo}, the edge platform of the company's race-timing service (596 of 605 commits): RFID readers, field units, Go services and the integration with the cloud API.
		\item Designed an \textbf{offline-first, store-and-forward architecture}: the reader publishes every tag to \textbf{RabbitMQ}, Go consumers persist reads across \textbf{sharded SQLite files} to avoid the single-writer bottleneck, and separate services sync event data and upload results when the network returns.
		\item Replaced the first Python/Flask monolith with small services (MyReader, Envio, Receba, Reenvio, Validate) and moved storage from MySQL to SQLite for reliability in the field.
		\item Rewrote the platform over one weekend to stream raw reads over \textbf{MQTT} and process them asynchronously in the cloud, with a parser and race logic on the server. Events ran without incidents from then on.
		\item Wrote a C++ driver for the Chafon RFID module with Python bindings through the CPython C API.
		\item Built the field unit's control panel: a Go daemon and Arduino firmware with a checksummed serial protocol, signal and battery monitoring, confirmations for destructive actions and power-down sleep.
		\item Ran production services (APIs, legacy systems, internal tools) on a VPS with rootless Podman Quadlets behind nginx, and wrote the install, update and supervision scripts for the field units.
		\item Built \textbf{PvAPI}, the Go backend of the company's PvAuto marketplace, replacing an earlier Laravel version: transactional checkout with row-level locking, per-shop earnings and withdrawals, partner API tokens and 19 HTTP test suites.
		\item Was the technical reference for the platform: onboarded the engineer who took over the services, setting up access and servers and walking him through the code.
	\end{itemize}

	%=== FEATURED PROJECTS ===
	\sectiontitle{Featured Projects}
	\begin{smallprojects}
		\textbf{MonoTempo} (Go, Python, C++, RabbitMQ, SQLite, Docker) -- \href{https://github.com/MyTempoESP/MonoTempo}{github.com/MyTempoESP/MonoTempo} \\
		Edge platform for RFID race timing, documented with architecture, design decisions and their trade-offs: offline first, a broker between reader and consumers, store everything and decide later, named pipes as the host control plane, and crash-and-restart instead of recovery logic.

		\textbf{PvAPI} (Go, Gin, pgx, PostgreSQL) -- \href{https://github.com/MyTempoESP/Pvapi}{github.com/MyTempoESP/Pvapi} \\
		Backend for a consignment marketplace, rewritten from Laravel to Go (489 commits, Mar--Jun 2026). Checkout runs in one transaction with row-level locks on the client, package and shop rows; money uses \texttt{decimal} and per-shop values are frozen at write time. Checkout idempotency went through three documented designs, from a client-supplied key to a content hash to a per-client rate limit after a production incident. Partner API with SHA-256-hashed tokens, per-entity authorization gates, 53 migrations and 19 black-box HTTP test suites against a real PostgreSQL. From May 2026 I ran it with coding agents: 64 written specs, guardrails in \texttt{AGENTS.md}, and custom skills for tests and diagrams.

		\textbf{A7X} (Lua, OpenResty/Lapis, MySQL) \\
		Race-results backend with 13 models, role-based access for organizers, judges and coaching groups, and reports. Built and deployed in three weeks.

		\textbf{UNIVC Analytics} (REST API, SQL, PHP) \\
		Enrollment analytics for the university. I built the student-data extraction API on top of the ERP database: a client creates an export, polls its status and downloads the CSV when it is ready. It replaced a scraper that logged into the JSF/RichFaces portal and downloaded Excel reports; each collection now takes a few seconds, so the dashboard syncs every hour and skips writes when a fingerprint of the reports hasn't changed.
	\end{smallprojects}

	%=== LANGUAGES ===
	\sectiontitle{Languages}
	Portuguese (Native) \quad | \quad English (Advanced)

	%=== EDUCATION ===
	\sectiontitle{Education}

	\textbf{Bachelor's in Computer Engineering} (in progress) \hfill Centro Universitário Vale do Cricaré \\
	Jul 2026 -- Jul 2030 (expected)

	\textbf{Technical High School, IT Support} \hfill CEEMTI Prof. Fernando Duarte Rabelo \\
	Jan 2021 -- Jan 2024

	\textbf{LiFT Scholarship recipient} \hfill The Linux Foundation \\
	Currently preparing for the Certified Kubernetes Administrator (CKA) certification

\end{document}
